\documentclass[UTF8,fontset=fandol,11pt]{ctexart}
\usepackage[a4paper,margin=24mm]{geometry}
\usepackage{amsmath,amssymb,tabularx,array,booktabs}
\usepackage[hidelinks]{hyperref}
\setlength{\parindent}{2em}
\setlength{\parskip}{0.3em}
\emergencystretch=2em
\title{六次势复标量孤子的固定荷敏感度与约束伴随交叉验证}
\author{}
\date{2026年9月27日}
\begin{document}
\maketitle

Fixed Charge Sensitivities and Constrained Adjoint Cross Validation for Complex Scalar Solitons with a Sextic Potential

OpenUFT · 2026年9月27日

\section*{摘要}

针对复标量孤子计算中作用量符号、径向伴随权重、固定荷约束及参数扫描不一致的问题，本文在有下界的六次势规范复标量模型中建立统一的驻点求导框架。通过消去静电场，推导光滑固定荷分支上的能量全微分，明确区分固定 Noether 荷和固定电荷的耦合导数。随后对球对称背景方程、频率、累计荷和渐近边界条件共同求导，实际计算电磁耦合、总荷及三个势能参数的敏感度。在选定参考分支上，完整能量导数积分与变分恒等式的最大绝对偏差约为3.83×10⁻¹¹，重新求解背景的中心差分表现出二阶收敛。另在中性固定荷模型中推导约束线性化的秩一修正，并以三档空间网格比较直接敏感度、伴随梯度及背景有限差分。漏掉固定荷修正时，示例导数由约+28.2312变为−737.9208，说明正确的局部径向算子不足以保证约束梯度正确。结果支持给定模型、扇区和分支上的变分与数值求导一致性；尚未证明带电全谱稳定性、真实动态辐射或现实粒子的对应关系。

\textbf{关键词}：Q-ball；六次势；Noether 荷；参数敏感度；约束伴随；电磁自能；数值验证

\section*{Abstract}

We examine parameter derivatives of stationary complex scalar solitons with a bounded sextic potential. Eliminating the electrostatic field yields an energy differential along smooth constrained stationary branches and distinguishes derivatives at fixed Noether charge from those at fixed electric charge. We differentiate the coupled spherical background equations, charge constraint, frequency and asymptotic boundary conditions, and compute sensitivities to the gauge coupling, total charge and three potential parameters. At a selected reference solution, energy derivatives obtained by integrating the complete field response agree with the variational identities within an absolute difference of approximately 3.83×10⁻¹¹. Central finite differences obtained by resolving the backgrounds exhibit second order convergence. For a neutral branch, we derive the rank one correction associated with fixed charge and compare direct sensitivities, adjoint gradients and background finite differences on three radial grids. Omitting this correction changes a representative derivative from approximately +28.2312 to −737.9208. These results support numerical and variational consistency within the specified model and branch. They do not establish stability of the full gauged system, dynamical radiation response, or an identification with an observed particle.

\textbf{Keywords}: Q-ball; sextic potential; Noether charge; sensitivity analysis; constrained adjoint; electrostatic energy.

\section*{1 引言}

统一场研究需要把几何假设、作用量、守恒量、目标解和观测规则连接起来。对于以局域场构型解释粒子的方案，首先需要一个自洽的场模型以及可复核的背景和导数。若场方程与能量泛函采用不同符号，或者参数优化时遗漏守恒约束，那么优化器的收敛不能形成有效物理证据。

Q-ball 是带全局相位和守恒荷的定态场构型。其稳定性、旋转结构与规范耦合需要分别处理；相关研究见文献[1,2]。本文不尝试由一个球对称经典背景同时解释电子自旋、强弱相互作用和引力，而集中研究这些更大目标所依赖的一段基础计算链。

本文的工作包括：统一修复模型与量纲；推导静电能的约束参数变分；实际求解五个参数的背景敏感度；通过背景重求解核对有限差分；在中性模型中构造完整固定荷伴随并进行空间收敛验证。采用的变分和伴随方法属于标准工具，本文不主张其方法本身为新理论。贡献在于对本候选模型执行相互独立的求导路径，并展示遗漏约束所导致的具体反例。

\section*{2 模型与适用范围}

采用 \(\hbar=c=1\) 和度规 \((-+++ )\)，作用量密度为

\[
\mathcal L=-(D_\mu\psi)^*D^\mu\psi-U(s)-\frac14F_{\mu\nu}F^{\mu\nu},
\quad s=|\psi|^2,
\tag{1}
\]
\[
D_\mu=\partial_\mu+ieA_\mu,
\quad U(s)=\mu^2s-\lambda s^2+g_6s^3.
\tag{2}
\]

要求 \(\lambda,g_6>0\) 且 \(\mu^2>\lambda^2/(4g_6)\)。于是 \(U(s)>0\) 对所有 s>0 成立。质量维数为 \([\psi]=[A_\mu]=1\)、\([\mu^2]=2\)、\([g_6]=-2\)，而 e 和 λ 无量纲。六次项可作为有效场论相互作用使用，其量子适用范围须另行指定截断尺度[3]。

Noether 电流及电流为

\[
j_N^\mu=-i[\psi^*D^\mu\psi-(D^\mu\psi)^*\psi],
\qquad J^\mu=e j_N^\mu.
\tag{3}
\]

取球对称静电定态 ansatz

\[
\psi=f(r)e^{-i\omega t},\quad a=eA_0,\quad W=\omega-a.
\tag{4}
\]

场方程化为

\[
f''+2f'/r=(\mu^2-W^2)f-2\lambda f^3+3g_6f^5,
\tag{5}
\]
\[
a''+2a'/r=-2e^2Wf^2,
\quad q'=8\pi r^2Wf^2.
\tag{6}
\]

这里 q 为累计 Noether 荷，\(q(\infty)=Q_N\)，电荷为 \(Q_{\rm el}=eQ_N\)。本文只研究指定球对称静电分支；其结果不覆盖完整非球对称动力学。

\section*{3 能量全微分的推导}

令 n 为 Noether 荷密度，G 为全空间静电 Green 函数，

\[
G(\mathbf x)=\frac1{4\pi|\mathbf x|},\quad a=e^2G*n.
\tag{7}
\]

消去静电势后的能量为

\[
E[f,n;e]=\int\left[|\nabla f|^2+U(f^2)+\frac{n^2}{4f^2}\right]d^3x
+\frac{e^2}{2}\iint n(\mathbf x)G(\mathbf x-\mathbf y)n(\mathbf y)d^3x d^3y.
\tag{8}
\]

对 \(E-\omega\int n\) 的 n 变分得到

\[
\frac{n}{2f^2}+a=\omega,
\qquad n=2Wf^2.
\tag{9}
\]

电磁能可以写成

\[
E_{\rm EM}=\frac12\int na\,d^3x
=\frac1{2e^2}\int|\nabla a|^2d^3x.
\tag{10}
\]

沿光滑驻点分支，总微分中的场变分满足 \(E_y\,dy=\omega\,dQ_N\)。在其余项中保留参数显式变化，得到

\[
\boxed{dE=\omega dQ_N+\frac{2E_{\rm EM}}e de
+I_2d\mu^2-I_4d\lambda+I_6dg_6},
\quad I_j=\int f^j d^3x.
\tag{11}
\]

该恒等式成立于光滑约束驻点，不要求它是全局极小值。e=0 处取连续极限。式(11)各项的质量维数均为1。

如果实际保持电荷 \(Q_{\rm el}\) 不变，则 \(dQ_N/de=-Q_N/e\)，相应导数为

\[
\left(\frac{\partial E}{\partial e}\right)_{Q_{\rm el}}
=\frac{2E_{\rm EM}-\omega Q_N}{e}.
\tag{12}
\]

本文的数值敏感度保持 Noether 荷不变，不能与式(12)混用。

\section*{4 背景和敏感度边值问题}

\subsection*{4.1 边界条件与尾能}

中心正则条件为 \(f'(0)=a'(0)=q(0)=0\)。有限端点 R 采用

\[
f'(R)+\beta f(R)=0,\quad a'(R)+a(R)/R=0,\quad q(R)=Q_N,
\tag{13}
\]
\[
\kappa=\sqrt{\mu^2-\omega^2},\quad
\gamma=\frac{e^2Q_N}{4\pi},\quad
\beta=\kappa+\frac{1+\omega\gamma/\kappa}{R}.
\tag{14}
\]

β 是指数局域尾部的领先阶 Coulomb 修正，不是精确有限域条件。结果筛选要求 \(0<\omega<\mu\) 和非零中心振幅。外部标量荷可忽略时，端点外静电能为

\[
E_{\rm out}=2\pi R^3a'(R)^2/e^2.
\tag{15}
\]

计算能量及导数时均包含该项，并通过扩大区域检查端点影响。

\subsection*{4.2 完整敏感度方程}

对 \(\eta\in\{e,Q_N,\mu^2,\lambda,g_6\}\)，令
\(h=\partial_\eta f\)、\(z=\partial_\eta a\)、\(v=\partial_\eta q\)、\(w_\eta=\partial_\eta\omega\)。定义
\(L=\mu^2-W^2-6\lambda f^2+15g_6f^4\)。逐项求导给出

\[
h''+2h'/r=Lh+2Wfz-2Wfw_\eta
+f\delta_{\eta,\mu^2}-2f^3\delta_{\eta,\lambda}+3f^5\delta_{\eta,g_6},
\tag{16}
\]
\[
z''+2z'/r=-4e^2Wfh+2e^2f^2z-2e^2f^2w_\eta
-4eWf^2\delta_{\eta,e},
\tag{17}
\]
\[
v'=16\pi r^2Wfh+8\pi r^2f^2(w_\eta-z).
\tag{18}
\]

中心条件为 \(h'(0)=z'(0)=v(0)=0\)，端点条件为

\[
h'(R)+\beta h(R)+f(R)(\beta_\omega w_\eta+\beta_\eta)=0,
\]
\[
z'(R)+z(R)/R=0,\qquad v(R)=\delta_{\eta,Q_N}.
\tag{19}
\]

需要使用的显式偏导为

\[
\beta_\omega=-\omega/\kappa+\gamma\mu^2/(R\kappa^3),
\quad\beta_e=2\omega\gamma/(e\kappa R),
\]
\[
\beta_Q=\omega e^2/(4\pi\kappa R),
\quad\beta_{\mu^2}=1/(2\kappa)-\omega\gamma/(2R\kappa^3),
\quad\beta_\lambda=\beta_{g_6}=0.
\tag{20}
\]

只对局部 ODE 求导而固定 ω 或忽略端点参数依赖，会改变原本的约束问题。

\section*{5 带电敏感度的数值结果}

\subsection*{5.1 基准与方法}

采用固定参考质量单位，设 \(\mu^2=\lambda=g_6=1\)、e=0.05、\(Q_N=279.1642942626188\)。改变参数时不重新定义参考单位。背景及敏感度采用含中心奇点矩阵的自适应边值求解，积分采用48001点 Simpson 求积。

R=100 的结果为

\[
\omega=0.960271793738176,\quad f(0)=0.712945292806638,
\]
\[
E=277.437311280806,\qquad E_{\rm EM}=1.65444352855416.
\tag{21}
\]

表1的能量导数由完整场响应积分得到，再与式(11)比较。

\textbf{表1 带电分支的参数导数}

\begin{center}
\small
\begin{tabularx}{\textwidth}{>{
aggedrightrraybackslash}X>{
aggedrightrraybackslash}X>{
aggedrightrraybackslash}X>{
aggedrightrraybackslash}X}
\toprule
参数 & \(\partial_\eta\omega\) & \(\partial_\eta E\) & 与式(11)的绝对偏差 \\
\midrule
e & 0.460174269398 & 66.177741142128 & \(3.83\times10^{-11}\) \\
\(Q_N\) & −0.000100266856251 & 0.960271793738 & \(2.18\times10^{-13}\) \\
\(\mu^2\) & 0.548965500136 & 147.174839367610 & \(3.15\times10^{-11}\) \\
λ & −0.177154489432 & −27.930524664504 & \(3.58\times10^{-11}\) \\
\(g_6\) & 0.068829603267 & 8.456183727213 & \(8.46\times10^{-13}\) \\
\bottomrule
\end{tabularx}
\end{center}

背景 Gauss 边界荷偏差约1.51×10⁻¹⁰。表中显示位数用于复核；误差指标不是严格区间认证，也不是现实观测精度。

\subsection*{5.2 参数差分与区域细化}

每次参数扰动都重新求解完整背景，采用中心差分。电磁耦合步长0.001、0.0005、0.00025时，能量差分与敏感度积分的绝对偏差依次约7.00×10⁻⁴、1.75×10⁻⁴、4.38×10⁻⁵。总荷步长1、0.5、0.25时对应偏差约1.37×10⁻⁷、3.41×10⁻⁸、8.53×10⁻⁹。其余势能参数也表现出预期的二阶差分规律。

将背景区域扩大到 R=140，背景容差由10⁻⁹收紧为10⁻¹⁰，得到 E=277.437311280804 和 \(\partial_eE=66.177741142111\)。两次能量相差约2.0×10⁻¹²，导数相差约1.8×10⁻¹¹。这是选定基准上的区域一致性检查，不能单独保证所有参数区域都同样准确。

\section*{6 中性固定荷伴随交叉验证}

\subsection*{6.1 约束线性化}

在 e=0 的中性分支，\(Q_N=2\omega I\)，消去频率后的能量是

\[
E_Q[f]=\int(|\nabla f|^2+U)d^3x+\frac{Q_N^2}{4I}.
\tag{22}
\]

驻点方程的约束线性化为

\[
K=L_++\frac{4\omega^2}{I}|f\rangle\langle f|,
\quad L_+=-\Delta+\mu^2-\omega^2-6\lambda f^2+15g_6f^4.
\tag{23}
\]

采用内积 \(\langle u,v\rangle=4\pi\int r^2uvdr\)，并在使边界形式消失的定义域上使用自伴性。对 J=I，敏感度及伴随满足

\[
Kf_\lambda=2f^3,\qquad Kz=2f,
\]
\[
\frac{dI}{d\lambda}=\langle2f,f_\lambda\rangle
=\langle z,2f^3\rangle.
\tag{24}
\]

\subsection*{6.2 离散实现与结果}

基准为 R=80、\(\omega=0.948683298051671\)、\(I=147.132501877045\)，总荷及参考势与前述一致。使用 u=rf 消去径向权重，采用二阶中心差分与齐次端点近似，并保留秩一约束项。

\textbf{表2 中性分支的空间收敛与伴随一致性}

\begin{center}
\small
\begin{tabularx}{\textwidth}{>{
aggedrightrraybackslash}X>{
aggedrightrraybackslash}X>{
aggedrightrraybackslash}X>{
aggedrightrraybackslash}X}
\toprule
网格步长 & 直接导数 & 伴随导数 & 二者绝对偏差 \\
\midrule
0.08 & 28.232524786890 & 28.232524786892 & \(1.94\times10^{-12}\) \\
0.04 & 28.231521327310 & 28.231521327304 & \(6.92\times10^{-12}\) \\
0.02 & 28.231270612587 & 28.231270612623 & \(3.62\times10^{-11}\) \\
\bottomrule
\end{tabularx}
\end{center}

每档直接与伴随使用同一离散矩阵，故其一致性验证转置关系，不能独自证明连续误差很小。为进一步检查，参数步长0.001、0.0005、0.00025的背景重求解差分给出28.231208861811、28.231192506240、28.231188417237。空间结果与最细参数差分的偏差约按四倍缩小，支持二阶收敛。

本离散测试以连续背景采样构建线性化矩阵，并非一个完整离散优化程序的端到端梯度认证。带电背景敏感度与本节中性伴随也是两个不同测试，不能将它们合并宣称完整带电伴随已经执行。

\subsection*{6.3 遗漏约束的反例}

在最细空间网格删除式(23)中的秩一项，得到 \(dI/d\lambda\approx-737.920785710305\)。这与固定荷导数在符号和数量级上均不同。其原因是固定频率与固定荷对应不同扰动空间，而不只是数值误差或精度不足。

\section*{7 讨论}

\subsection*{7.1 能量源项必须从同一泛函求导}

对 \(E_0=4\pi\int r^2[\tfrac12f'^2+V_1f^4/4-V_2f^6/6]dr\)，加权泛函导数为 \(-\Delta_rf+V_1f^3-V_2f^5\)。在与其一致的静态驻点上，这一源项为零；只有使用相反非线性符号的背景时，才会得到两倍势能导数。不能把混用模型产生的因子作为通用伴随修复。

此外，正 V₁ 条件下的三维质量为零静态方程，在正则有限积分及消失边界项假设下，积分恒等式与 Pohozaev 恒等式迫使 \(\int f^4=0\)。本文的定态修复模型因此是另选候选，不是原静态模型的数值确认。

\subsection*{7.2 求导验证与稳定性不同}

式(11)及其数值核对验证驻点与参数变分的一致性。\(\partial\omega/\partial Q_N<0\) 是分支信息，不能独自证明带电扰动稳定。完整稳定性还需处理 Gauss 约束、规范零模、耦合动力学谱和非球对称扇区。

球对称静电标量背景也没有产生空间电流磁矩或电子内禀自旋。因此，把经典背景能量指定为电子质量，再拟合电荷，不能完成现实粒子的识别。相关量子态与实际观测映射须另行建立。

\subsection*{7.3 精度与参数扫描}

本文计算采用双精度。提高某个内部积分库的任意精度设置，不能提升外部双精度优化器的参数精度[4,5]。每个扫描参数点应重新求解背景并检查可行性；沿旧轮廓计算新参数积分，不代表新参数的驻点。

参数数量大于观测数量不保证存在平滑简并流形；需要观测映射的秩与局部正则性。附加约束也只有在相交和横截等条件下改变局部维数。本文没有以参数数量代替可识别性证明。

\section*{8 结论}

本文在指定六次势规范复标量模型中完成固定 Noether 荷参数响应与能量变分的交叉验证，并在中性模型中验证包含秩一约束的伴随梯度。解析能量全微分、完整背景敏感度积分、重新求解的有限差分和空间细化结果相互一致。遗漏固定荷响应会产生显著错误，表明约束与边界求导是计算链的必要部分。

结论适用于所研究模型、球对称静电扇区和选定光滑分支。后续工作应优先构造完整带电约束扰动谱、真实动态响应及量子观测定义。本文不把这些数值与变分结果表述为全物理统一或电子性质的独立预言。

\section*{数据与复现说明}

数值数据保存于项目的 \texttt{V3\_4\_fixed\_charge\_sensitivity.json} 和 \texttt{V3\_5\_adjoint\_crosschecks.json}，中性轮廓保存于 \texttt{V3\_Qball\_profile.csv}。使用 Python 3.8.8、NumPy 1.24.3、SciPy 1.10.1；背景最大节点数分别为60000和50000。实际步长、容差、能量尾部及方法见第4–6节和数据文件。

背景代码在研究会话中执行，本文随稿提供计算定义和原始结果，不宣称已提供独立可运行的复现软件包。投稿前应补齐完整复现程序、来源版本与自动生成表格的流程，并确认作者署名、数据许可和目标期刊要求。

\section*{参考文献}

[1] A. G. Panin and M. N. Smolyakov. Classical behaviour of Q-balls in the Wick–Cutkosky model. European Physical Journal C 79, 150 (2019). \href{https://link.springer.com/article/10.1140/epjc/s10052-019-6638-2}{DOI 10.1140/epjc/s10052-019-6638-2}.

[2] Mikhail S. Volkov and Erik Woehnert. Spinning Q-Balls. arXiv:hep-th/0205157 (2002). \href{https://arxiv.org/abs/hep-th/0205157}{预印本}.

[3] David B. Kaplan. Five lectures on effective field theory. arXiv:nucl-th/0510023 (2005). \href{https://arxiv.org/abs/nucl-th/0510023}{讲义}.

[4] SciPy. minimize method SLSQP. \href{https://docs.scipy.org/doc/scipy/reference/optimize.minimize-slsqp.html}{官方文档}.

[5] mpmath. Precision and representation issues. \href{https://www.mpmath.org/doc/current/technical.html}{官方文档}.

[6] OpenUFT. 固定荷敏感度求导与能量交叉验证, 2026.

[7] OpenUFT. 伴随审查复核与约束梯度融合, 2026.
\end{document}